% 2026-07-13: \glsfirst anchors moved out of the section/subsection titles (underlined
% glossary links don't belong in headings); the anchors now sit at the first body uses.
\section{Algebra: what is a learned feature?}
\label{sec:interp}

\emph{Mechanistic \glsfirst{interpretability}{interpretability}} tries to reverse-engineer a trained neural network into something a human can read. Much of its core is linear algebra and representation theory: how a network stores more concepts than it has dimensions, and what it means for a concept to be a \emph{direction} one can read or steer.

\subsection{Superposition: packing features into directions}

How does a network with a few hundred dimensions in each layer represent
the many thousands of distinct concepts (``\glsfirst{feature}{features}'')
it seems to know?
A \glsfirst{transformer}{transformer} reads text as a sequence of
\glsfirst{token}{tokens}, words and word-fragments; at any fixed layer,
the computation has the form
\[
\text{text}
\longrightarrow
\text{tokens}
\longrightarrow
\underbrace{(h_1,\ldots,h_T)\in(\mathbb{R}^n)^T}
            _{\text{activations at one layer}}
\longrightarrow
\text{next-token log-probabilities}.
\]
For a text of $T$ tokens, each token position carries a vector
$h_j\in\mathbb{R}^n$.
The copy of $\mathbb{R}^n$ in which these vectors live is the network's
\glsfirst{activation-space}{activation space}.
A leading hypothesis, made vivid by Elhage et al.~\cite{elhage2022},
is \glsfirst{superposition}{superposition}: when features are sparse,
a network can store more of them than it has dimensions, as directions in
activation space that are \emph{not} orthogonal but only nearly so,
tolerating the resulting interference because at any moment only a few
features are active (Figure~\ref{fig:superposition}).
In a small \glsfirst{relu}{ReLU} model one can watch the geometry organize
itself into regular polytopes---antipodal pairs, triangles, pentagons---as
a function of how sparse and how important the features are; it is,
recognizably, a Thomson problem (where do mutually repelling charges
on a sphere come to rest?)~\cite{saff1997}.

\begin{figure}[ht]
\centering
\begin{tikzpicture}[>=Stealth, scale=1.15]
  \begin{scope}
    \node[font=\small] at (0,1.65) {dictionary directions};
    \draw[->, gray!60] (-1.35,0)--(1.35,0);
    \draw[->, gray!60] (0,-1.35)--(0,1.35);
    \foreach \a/\i in {90/1, 162/2, 234/3, 306/4, 18/5} {
      \draw[->, thick, gray!80!black] (0,0) -- (\a:1);
      \node at (\a:1.2) {\footnotesize $f_{\i}$};
    }
    \draw[->, very thick, blue!60!black] (0,0) -- (90:1);
    \draw[->, very thick, blue!60!black] (0,0) -- (306:1);
  \end{scope}
  \begin{scope}[xshift=4.7cm]
    \node[font=\small] at (0,1.65) {one sparse activation};
    \draw[->, gray!60] (-1.35,0)--(1.35,0);
    \draw[->, gray!60] (0,-1.35)--(0,1.35);
    \foreach \a in {90,162,234,306,18} {
      \draw[->, thin, gray!65] (0,0) -- (\a:1);
    }
    \coordinate (a) at (90:1);
    \coordinate (b) at (306:1);
    \coordinate (sum) at ($(a)+(b)$);
    \draw[->, very thick, blue!60!black] (0,0) -- (a) node[above] {\footnotesize $f_1$};
    \draw[->, very thick, blue!60!black] (a) -- (sum) node[midway, right] {\footnotesize $f_4$};
    \draw[->, very thick, red!65!black] (0,0) -- (sum) node[right] {\footnotesize $y=f_1+f_4$};
    \draw[dashed, gray!70] (sum) -- (b);
  \end{scope}
  \draw[->, thick] (1.6,0) -- (3.0,0)
    node[midway, above, font=\scriptsize, align=center] {few active\\at once};
\end{tikzpicture}
\caption{Superposition: five feature directions packed into a two-dimensional space as a regular pentagon. No two are orthogonal, but if only a few features are active at once, the activation is a sparse sum such as $y=f_1+f_4$ and the active directions can still be recovered despite interference.}
\label{fig:superposition}
\end{figure}

The nearest classical analogue is \emph{compressed sensing}~\cite{candes2006,donoho2006}. Picture the features active on a given input as a sparse vector $x\in\mathbb{R}^m$ (only a few of the $m$ possible features fire at once); the network stores it as $y=\Phi x$ in the $n\ll m$ dimensions of its activation space---a linear \emph{measurement}---and reading a feature back out is the \emph{recovery} of $x$ from $y$. Two classical facts explain how this can work. First, there is room for the directions: for every $\varepsilon\in(0,1)$, the space $\mathbb{R}^n$ contains $e^{\Omega(\varepsilon^2 n)}$ unit vectors whose pairwise inner products are at most $\varepsilon$ in absolute value, so the number of $\varepsilon$-almost-orthogonal directions grows \emph{exponentially} in the dimension (a consequence of the Johnson--Lindenstrauss lemma~\cite{johnson1984}). And second, when only a few of those directions are active at once, the stored signal can be read back out. Say that $\Phi\in\mathbb{R}^{n\times m}$ is \emph{restricted-isometric of order $k$} with distortion $\delta$ if
\[
   (1-\delta)\|x\|^2\le\|\Phi x\|^2\le(1+\delta)\|x\|^2
\]
for every \emph{$k$-sparse} vector $x$, meaning one with at most $k$ nonzero entries.

\begin{theorem*}[Cand\`es~\cite{candes2008}; sharp constant, Cai--Zhang~\cite{caizhang2014}]
If $\Phi$ is restricted-isometric of order $2s$ with distortion $\delta<1/\sqrt{2}$, then every $s$-sparse $x\in\mathbb{R}^m$ is the \emph{unique} minimizer of
\[
   \min_{z}\ \|z\|_1 \quad\text{subject to}\quad \Phi z=\Phi x,
\]
and so is recovered \emph{exactly} by a convex program. A random $\Phi$, say with independent Gaussian entries suitably normalized, has the property with high probability once
\[
   n \;\gtrsim\; s\,\log(m/s).
\]
\end{theorem*}

\noindent That last scaling is the quantitative heart of the analogy: the dimension a network needs grows linearly with the number of features active \emph{at once}, and only logarithmically with the total number it can represent. A few hundred dimensions can carry many thousands of features, provided only a handful fire on any given input. For trained transformers this remains an analogy rather than a theorem; turning it into one is part of the invitation.

There is a catch.
Compressed sensing assumes the dictionary $\Phi$ is \emph{known}.
What if someone hands you a trained network and asks what features it
represents?
You can run inputs through the network and record activation vectors,
but the feature directions were not supplied with the
\glsfirst{weights}{weights} (the network's trained parameters).
Recovering the features is therefore a problem in \emph{dictionary learning}:
from the observed activation vectors alone, simultaneously infer a dictionary
$\Phi$ whose columns are feature directions and, for each activation $y$,
a sparse coefficient vector $x$ with $y\approx\Phi x$.
Both the dictionary $\Phi$ and the codes $x$ are unknown.

A \glsfirst{sparse-autoencoder}{sparse autoencoder} attempts this empirically.
It is trained to reconstruct activation vectors using a learned dictionary,
with an $\ell^1$ penalty encouraging each reconstruction to use only a few
dictionary directions.
The resulting directions are often more interpretable than the network's raw
neurons~\cite{bricken2023,cunningham2023}.
This raises an identifiability question: when does the sparse autoencoder
recover features already used by the network, rather than impose a new
coordinate system that merely makes the activations easier to describe?
Classical dictionary learning answers this question when the codes are
exactly sparse and their supports well spread: the dictionary is then
determined up to relabeling and rescaling of its
columns~\cite{aharon2006,hillar2015}, stably in the presence of
noise~\cite{gribonval2015,garfinkle2019}.
Network activations flout these hypotheses because their features co-occur,
which is where the first open problem below begins.

\subsection{Reading, steering, and learned representations}

What could it mean for a concept to be a direction in activation space?
Consider incomplete sentences such as ``The keys on the table \ldots'' and
pause the computation at the activation $h_T$ of the last token.
A linear functional $\ell$ \emph{reads} grammatical number if, for some
threshold $c$, the inequality $\ell(h_T)>c$ predicts that the subject is
plural.
A vector $v\in\mathbb{R}^n$ \emph{steers} toward the plural if replacing
$h_T$ by $h_T+v$ raises the log-probability of ``are'' relative to ``is.''
Thus the same example produces a readout covector $\ell$ and a steering
vector $v$.
Relating them requires an inner product.

Park, Choe, and Veitch~\cite{park2023} formalize these two experiments for
next-token prediction.
They ask whether activation space carries a \emph{causal inner product}
in which independently varying concepts are orthogonal and whose Riesz map
sends readout covectors to steering vectors.
They estimate this geometry from the unembedding matrix
(the model's final linear map, from activation space to next-token
scores).

Arditi et al.~\cite{arditi2024} show that in several open chat models, refusal is mediated by a single direction in the \glsfirst{residual-stream}{residual stream} (the running activation vector that a transformer's layers read from and write to): this direction both predicts whether the model will refuse and causally controls refusal when added or removed. A lone direction serving as both readout and control is evidence for the linear representation hypothesis, though in a different representation space from the unembedding geometry above.

Refusal and its cousins are semantic properties, with no ground truth for what the network \emph{should} compute. Arithmetic is a cleaner laboratory: train a network on a group operation, and the algorithm it ought to discover is known in advance. Representation theory appears in trained arithmetic circuits. A one-layer transformer trained to add integers modulo $p$ learns a discrete-Fourier ``clock'' algorithm, with the characters of $\mathbb{Z}/p\mathbb{Z}$ appearing in its weights~\cite{nanda2023}. Chughtai, Chan, and Nanda~\cite{chughtai2023} take the first step beyond this abelian case. The ambient structure is the Artin--Wedderburn decomposition of the group algebra into matrix blocks,
\[
   \mathbb{C}[G]\;\cong\;\bigoplus_{\rho\in\widehat{G}} M_{d_\rho}(\mathbb{C}),
   \qquad \sum_{\rho\in\widehat{G}} d_\rho^{2}=\lvert G\rvert,
\]
one block per irreducible representation $\rho$. For a nonabelian group the higher-dimensional blocks are genuine matrix representations, and trained networks compute the product through a sparse subset of them. The algorithm's correctness is a representation-theory theorem; that networks implement it is empirical. Which irreducible representations training selects is still something we discover after the fact.

Sometimes the emergence of this structure is a theorem rather than an observation. Marchetti et al.~\cite{marchetti2024} prove that for architectures whose weights are pinned down by the function they compute (up to a unitary change of basis in each unit), invariance under a finite group $G$ forces each unit's weights to be an irreducible unitary representation of $G$; when the weights are moreover orthonormal, together they form the Fourier transform on $G$, and the group itself can be read off the trained weights.

What makes this a safety topic is \emph{selection}: a trained network can discover a crisp algebraic algorithm, hide it in weights, and use only part of the mathematically available structure. If we want to predict a system's behavior off-distribution, it matters which algorithm training found, not only that the training \glsfirst{loss}{loss} is low.

One caveat carries across this whole enterprise: a direction that \emph{predicts} a feature is only a \emph{correlate}. Showing that the network actually \emph{uses} it takes an intervention---ablating the direction, or patching it to the value it would take on another input, and watching the behavior change---not a \glsfirst{probe}{probe} that merely reads it off.

\subsection{Open problems}

% pass 2: moved here from the end of the section, so the section ends on the last open problem rather than on a bibliography pointer
\noindent A broad, structured catalogue of open problems in this area---several with a precise mathematical core---is collected by Sharkey et al.~\cite{sharkey2025}.

\begin{openproblem}[The geometry and identifiability of superposition]
Suppose activations have the form $y=\Phi x+\xi$, where the columns $v_1,\dots,v_m$ of $\Phi$ are unit feature directions, $\xi$ is noise, and the sparse codes $x$ have correlated supports. Estimate the dictionary the way a sparse autoencoder does: minimize reconstruction error plus an $\ell^1$ penalty. Determine, in terms of the coherence $\mu=\max_{i\neq j}\lvert\langle v_i,v_j\rangle\rvert$, the sparsity, the sample size, and the support correlations, when every minimizer recovers the true directions up to permutation and scaling, and when it instead \emph{merges} co-occurring directions into one. Already for two features the answer is unknown: a nested two-event model provably forces merging for every sufficiently small positive penalty, while solo firing restores recovery in the zero-penalty constrained limit; locating the positive-penalty phase boundary is the place to start.
\end{openproblem}

\begin{openproblem}[Training for interpretability]
Post-hoc interpretability asks whether a trained network's features can be recovered after the fact. A complementary problem is to train networks whose features are easier to recover in the first place. In the ReLU toy model of Elhage et al.~\cite{elhage2022} the features are known by construction, so the trade-off can be made exact. Determine the interference--performance frontier: for features appearing sparsely and independently, among weight matrices whose coherence (worst-case interference between stored features) is at most $\mu$, how small can the task loss be, as a function of $\mu$, the sparsity, the number of features, and the number of neurons? A first case: prove or refute that penalizing the \emph{average} interference during training lowers the \emph{worst-case} coherence of the minimizer, compared with training on the task loss alone.
\end{openproblem}

\begin{openproblem}[Representation theory of learned circuits]
Networks trained on group multiplication learn representation-theoretic algorithms~\cite{nanda2023,chughtai2023}, but \emph{which} irreducible representations they use varies from one random seed to the next. Fix a finite group $G$, a one-hidden-layer architecture of fixed width, Gaussian initialization, and gradient flow on the cross-entropy loss (the mean negative log-probability assigned to the correct product) with \emph{weight decay}, an added penalty proportional to the squared norm of the weights. The irreducibles visible in the trained network's outputs then form a random subset of $\widehat{G}$; determine its distribution---which irreducibles are learned, with what probability, as a function of the width and the decay strength. The tables of learned representations in~\cite{chughtai2023} are the data such a theorem would have to explain; the smallest case with genuine competition, $G=\mathbb{Z}/5\mathbb{Z}$ at width two, is open already with the decay switched off.
\end{openproblem}

% \noindent\emph{Entry points.} Linear algebra, compressed sensing, harmonic analysis, representation theory, and differential geometry.
